\section{Background and Motivation}
We selected this paper because our group is interested in exploring how clustering can support retrieval-augmented generation (RAG) for biomedical data, which is a growing research area. Single-cell transcriptomic data is a good fit for this direction because it is large, high-dimensional, and often difficult to interpret directly. The "scMINER" framework is especially suitable because it introduces a clustering approach based on mutual information that is designed to capture complex relationships among cells and identify meaningful cell populations. This gives us a strong clustering foundation for the course while also creating a natural next step where the important genes or markers from each cluster can be used to retrieve relevant biomedical literature through RAG. In this way, the project can remain focused on clustering while also extending the results toward a modern biomedical interpretation workflow.
